\chapter{Introducere}\label{ch:intro}
\pagestyle{fancy}

\section{Context și motivație}
În ultimii ani, aplicațiile de viziune artificială au devenit uzuale în industrie și cercetare, iar calitatea lor depinde direct de seturile de date etichetate. Procesul de adnotare vizuală implică trasarea manuală a obiectelor în imagini, fie prin dreptunghiuri de încadrare, fie prin poligoane sau măști. Pentru seturi mari de date, această muncă este consumatoare de timp și necesită atenție constantă pentru a menține consistența între adnotări. În practică, multe proiecte întârzie sau cresc în cost tocmai din cauza efortului de etichetare.

Motivația lucrării pornește din nevoia de a reduce timpul și efortul necesare adnotării, fără a pierde controlul uman asupra rezultatului final. O abordare semi-automată, în care un model de inteligență artificială propune contururi sau măști, poate accelera semnificativ procesul. Utilizatorul rămâne totuși responsabil de validare și corectare, ceea ce păstrează calitatea datelor. În plus, exportul adnotărilor în formate standard ajută la integrarea rapidă a datelor în fluxuri de antrenare.

Un alt aspect important este colaborarea între mai mulți utilizatori. În proiectele reale, adnotarea se face de obicei în echipă, iar lipsa sincronizării duce la duplicarea muncii și la inconsistențe. Printr-un mecanism de lucru colaborativ, cu actualizări în timp real și gestionarea sesiunilor, echipa poate lucra pe același proiect în mod coordonat. Astfel, se obține un proces mai rapid, mai controlat și mai ușor de urmărit.

\section{Domeniul și delimitarea temei}
Lucrarea se încadrează în domeniul adnotării vizuale pentru date de tip imagine, cu accent pe segmentare și detecție de obiecte. Sunt vizate adnotări realizate prin dreptunghiuri, poligoane și măști, deoarece acestea acoperă cele mai folosite reprezentări în seturile de date moderne.

Domeniul ales include și administrarea etichetelor și a adnotărilor, astfel încât acestea să poată fi organizate și verificate ușor.

Datele rezultate sunt salvate într-un format intern și pot fi exportate în formate standard, precum COCO (Common Objects in Context) și YOLO (You Only Look Once). Astfel, seturile de date pot fi folosite ușor în fluxuri de antrenare, fără pași suplimentari de conversie. Structura de export este gândită pentru a păstra separarea între tipurile de adnotări și pentru a face reutilizarea cât mai simplă.

Tema este delimitată la un instrument desktop de adnotare, care integrează metode semi-automate prin apelarea unor modele externe. Modelul nu este antrenat în cadrul lucrării, ci este folosit ca suport pentru propunerea contururilor. Pentru segmentare asistată este folosit un model de tip Mask2Former, iar pentru detecție și segmentare completă se folosește un model de tip YOLOE.

Nu sunt urmărite subiecte precum optimizarea arhitecturilor de rețele neuronale, compararea extensivă a modelelor sau dezvoltarea unui serviciu cloud pentru adnotare. Accentul este pe realizarea unui instrument practic, ușor de folosit, local, care reduce efortul de etichetare și permite colaborarea eficientă între membri ai aceleiași echipe.

\section{Soluția propusă și contribuții}
Soluția propusă este o aplicație desktop dezvoltată în PyQt6, care oferă un flux de lucru complet pentru adnotarea imaginilor. Utilizatorul poate deschide un folder de imagini, alege rapid fișierele din listă și edita adnotările direct pe canvas. Interfața este organizată în panouri funcționale, astfel încât instrumentele de adnotare, lista de elemente și proprietățile acestora să fie accesibile fără a întrerupe procesul de lucru.

Canvasul este componenta centrală a interfeței. Acesta permite zoom și pan, desenarea de dreptunghiuri, poligoane sau măști, precum și selecția și modificarea elementelor existente. În plus, sunt disponibile ajustări vizuale ale imaginii pentru a crește claritatea în zonele dificil de adnotat. Acest mod de lucru ajută la păstrarea preciziei și reduce efortul în sarcinile repetitive.

Gestionarea adnotărilor se face printr-un manager central, care păstrează toate elementele asociate imaginii curente și oferă operații de adăugare, modificare sau ștergere. Managerul emite notificări către interfață atunci când apar schimbări, ceea ce asigură sincronizarea dintre starea internă și elementele afișate. Fiecare adnotare are un identificator propriu, fapt util pentru editare, export și colaborare.

Aplicația include facilități de tip quality of life, menite să reducă erorile și timpul pierdut. La închiderea aplicației sau la schimbarea imaginii, utilizatorul este întrebat dacă dorește să salveze modificările. În plus, aplicația include chat în timp real și suport de tip chatbot local pentru asistență.

Componenta de inteligență artificială este integrată prin AutoSeg și AutoMask, care rulează modele externe ca subprocese asincrone. Rezultatele sunt transformate în adnotări editabile, iar utilizatorul le poate corecta imediat. Astfel, aplicația combină propunerea automată cu controlul manual, obținând un echilibru între viteză și acuratețe.

Contribuțiile principale constau în integrarea unui flux complet de adnotare cu suport pentru segmentare semi-automată, organizarea clară a interfeței și mecanismele de sincronizare în colaborare. Rezultatul este un instrument practic, adaptat nevoilor de etichetare din proiecte reale.

\section{Arhitectură colaborativă}
Backend-ul aplicației este separat pe servicii dedicate pentru autentificare, chat și colaborare pe adnotări. Un reverse proxy direcționează cererile către serviciul corespunzător, iar mesageria facilitează schimbul de informații între module. Această separare permite dezvoltarea și testarea componentelor în mod independent, păstrând totuși un flux unitar pentru utilizator.

Sincronizarea se bazează pe două tipuri de comunicație. Prin REST, clientul listează sesiunile, obține un snapshot al stării curente și poate încărca sau descărca imagini temporare. Prin WebSocket, clienții primesc evenimente în timp real, precum crearea, actualizarea sau ștergerea adnotărilor. Când o imagine este actualizată într-o sesiune, ceilalți utilizatori primesc un eveniment de tip „image available” și sincronizează automat datele.

Modelul de colaborare este centrat pe sesiuni, fiecare sesiune având o stare curentă și o versiune a adnotărilor. Astfel, modificările sunt propagate rapid, iar utilizatorii văd același conținut fără conflicte majore. Autentificarea prin JWT (JSON Web Token) limitează accesul doar la utilizatorii autorizați, iar token-ul este salvat local pentru a evita relogarea repetată.

